\documentclass[12pt,a4paper]{article}
\usepackage[utf8]{inputenc}
\usepackage[T1]{fontenc}
\usepackage[provide=*,main=portuguese]{babel}
\usepackage[a4paper,margin=2.5cm]{geometry}
\usepackage{graphicx}
\usepackage{float}
\usepackage{hyperref}
\usepackage{xcolor}
\usepackage{listings}
\usepackage{listingsutf8}
\usepackage{amssymb}
\usepackage{caption}
\renewcommand{\figurename}{Figura}
\captionsetup[figure]{labelfont=bf,labelsep=colon,justification=centering}
\lstset{
  basicstyle=\ttfamily\small,
  breaklines=true,
  columns=fullflexible,
  frame=single,
  framerule=0.4pt,
  rulecolor=\color{black!30},
  xleftmargin=0.8em,
  xrightmargin=0.8em,
  inputencoding=utf8,
  extendedchars=true,
  literate={á}{{\'a}}1 {à}{{\`a}}1 {â}{{\^a}}1 {ã}{{\~a}}1
           {Á}{{\'A}}1 {À}{{\`A}}1 {Â}{{\^A}}1 {Ã}{{\~A}}1
           {é}{{\'e}}1 {ê}{{\^e}}1 {É}{{\'E}}1 {Ê}{{\^E}}1
           {í}{{\'i}}1 {Í}{{\'I}}1
           {ó}{{\'o}}1 {ô}{{\^o}}1 {õ}{{\~o}}1
           {Ó}{{\'O}}1 {Ô}{{\^O}}1 {Õ}{{\~O}}1
           {ú}{{\'u}}1 {Ú}{{\'U}}1
           {ç}{{\c{c}}}1 {Ç}{{\c{C}}}1
           {º}{{\textordmasculine}}1 {ª}{{\textordfeminine}}1
           {✓}{{\checkmark}}1 {•}{{\textbullet}}1
           {π}{{$\pi$}}1 {ω}{{$\omega$}}1 {Ω}{{$\Omega$}}1 {ζ}{{$\zeta$}}1 {τ}{{$\tau$}}1
           {≈}{{$\approx$}}1
           {≤}{{$\le$}}1 {≥}{{$\ge$}}1
           {²}{{$^{2}$}}1 {³}{{$^{3}$}}1
           {±}{{$\pm$}}1
           {→}{{$\rightarrow$}}1 {✗}{{x}}1 {─}{{-}}1 {│}{{|}}1
           {ₙ}{{$_n$}}1
           {₀}{{$_0$}}1 {₁}{{$_1$}}1 {₂}{{$_2$}}1 {₃}{{$_3$}}1 {₄}{{$_4$}}1
           {₅}{{$_5$}}1 {₆}{{$_6$}}1 {₇}{{$_7$}}1 {₈}{{$_8$}}1 {₉}{{$_9$}}1,
}

\title{Trabalho 1: Desempenho de Sistemas de Controle e Projetos}
\author{Igor de Matos da Rosa \\ Matrícula: 20103930 \\ Universidade Federal de Santa Catarina \\ Disciplina: Controle Aplicado à Computação}
\date{25/04/2026}

\begin{document}
\maketitle
\thispagestyle{empty}

\clearpage
\section*{PC4.1}
\noindent\textit{Fonte:} 4.1.ipynb
\vspace{0.4cm}
\subsection*{Resultados (saída do código)}
\begin{lstlisting}
Máxima Ultrapassagem (MP): 35.09%

Valor final (y_ss): 1.20

Erro em regime estacionário (e_ss): -0.20
\end{lstlisting}
\subsection*{Gráficos}
\begin{figure}[H]
\centering
\includegraphics[width=0.95\linewidth]{figs/4\_1\_fig01.png}
\caption{PC4.1}
\end{figure}

\clearpage
\section*{PC4.2}
\noindent\textit{Fonte:} 4.2.ipynb
\vspace{0.4cm}
\subsection*{Resultados (saída do código)}
\begin{lstlisting}
Ganho DC da função de transferência: 0.2

Valor da saída em regime estacionário: 0.2000
Erro em regime estacionário (1 - saída): 0.8000

Como esperado, o erro é aproximadamente 0.8

==================================================
RESULTADO DO EXERCÍCIO PC4.2
==================================================
Função de transferência: G(s) = 4/(s² + 2s + 20)
Entrada: Degrau unitário
Valor desejado em regime estacionário: 1
Valor obtido em regime estacionário: 0.2
Erro em regime estacionário: 0.8
==================================================
\end{lstlisting}
\subsection*{Gráficos}
\begin{figure}[H]
\centering
\includegraphics[width=0.95\linewidth]{figs/4\_2\_fig01.png}
\caption{PC4.2}
\end{figure}
\begin{figure}[H]
\centering
\includegraphics[width=0.95\linewidth]{figs/4\_2\_fig02.png}
\caption{PC4.2}
\end{figure}

\clearpage
\section*{PC4.3}
\noindent\textit{Fonte:} 4.3.ipynb
\vspace{0.4cm}
\subsection*{Resultados (saída do código)}
\begin{lstlisting}
Funções de Transferência:
K = 10: T(s) = 50 / (s² + 15s + 10)
  Polos: [-14.30073525  -0.69926475]

K = 200: T(s) = 1000 / (s² + 15s + 200)
  Polos: [-7.5+11.98957881j -7.5-11.98957881j]

K = 500: T(s) = 2500 / (s² + 15s + 500)
  Polos: [-7.5+21.06537443j -7.5-21.06537443j]

=============================================================================
ANÁLISE DAS RESPOSTAS AO DEGRAU UNITÁRIO
=============================================================================

Função de Transferência em Malha Fechada: T(s) = 5K / (s² + 15s + K)

  K Ganho DC Máximo Sobressinal (%) Tempo de Acomodação (s) Erro em Regime Est.
 10   4.3548                   0.00                  2.8199             -3.3548
200   5.0000                  14.01                  0.4097             -4.0000
500   5.0000                  32.67                  0.4967             -4.0000
=============================================================================

ANÁLISE:
-----------------------------------------------------------------------------

K = 10:
  • Ganho DC (y_ss): 4.3548
  • Máximo Sobressinal: 0.00%
  • Tempo de Acomodação (±2%): 2.8199 s
  • Erro em Regime Estacionário: -3.3548

K = 200:
  • Ganho DC (y_ss): 5.0000
  • Máximo Sobressinal: 14.01%
  • Tempo de Acomodação (±2%): 0.4097 s
  • Erro em Regime Estacionário: -4.0000

K = 500:
  • Ganho DC (y_ss): 5.0000
  • Máximo Sobressinal: 32.67%
  • Tempo de Acomodação (±2%): 0.4967 s
  • Erro em Regime Estacionário: -4.0000

=============================================================================
ANÁLISE COMPLEMENTAR - CARACTERÍSTICAS DOS SISTEMAS
=============================================================================

K = 10:
  • Denominador: s² + 15s + 10
  • Frequência Natural (ωₙ): 3.1623 rad/s
  • Fator de Amortecimento (ζ): 2.3717
  • Sistema: Superamortecido (overdamped)
  • Constantes de tempo: τ₁ = 0.0699 s, τ₂ = 1.4301 s

K = 200:
  • Denominador: s² + 15s + 200
  • Frequência Natural (ωₙ): 14.1421 rad/s
  • Fator de Amortecimento (ζ): 0.5303
  • Sistema: Subamortecido (underdamped)
  • Frequência amortecida (ωd): 11.9896 rad/s
  • Frequência amortecida (f_d): 1.9082 Hz
  • Tempo de Pico (t_p): 0.2620 s

K = 500:
  • Denominador: s² + 15s + 500
  • Frequência Natural (ωₙ): 22.3607 rad/s
  • Fator de Amortecimento (ζ): 0.3354
  • Sistema: Subamortecido (underdamped)
  • Frequência amortecida (ωd): 21.0654 rad/s
  • Frequência amortecida (f_d): 3.3527 Hz
  • Tempo de Pico (t_p): 0.1491 s

=============================================================================
\end{lstlisting}
\subsection*{Gráficos}
\begin{figure}[H]
\centering
\includegraphics[width=0.95\linewidth]{figs/4\_3\_fig01.png}
\caption{PC4.3}
\end{figure}

\clearpage
\section*{PC4.5}
\noindent\textit{Fonte:} 4.5.ipynb
\vspace{0.4cm}
\subsection*{Resultados (saída do código)}
\begin{lstlisting}
=============================================================================

VERIFICAÇÃO FINAL COM k = 4.3081
=============================================================================

Parâmetros do Sistema:
  Função de Transferência: T(s) = 10/(s² + 4.3081s + 10)

Características da Resposta:
  • Ganho DC (y_ss):                0.997875
  • Máximo Sobressinal (Mp):        5.60%
  • Tempo de Pico (t_p):            1.3567 s
  • Tempo de Acomodação (±2%):     1.9315 s
  • Erro em Regime Estacionário:    0.002125

✓ CRITÉRIO SATISFEITO: 5.60% está dentro do intervalo [1%, 10%]

✓ ERRO EM REGIME ESTACIONÁRIO: 0.00212468 ≈ 0 (praticamente nulo)

=============================================================================

CONCLUSÃO:
-----------------------------------------------------------------------------
O valor k = 4.3081 foi selecionado para satisfazer:
  1. Máximo sobressinal entre 1% e 10%: Mp = 5.60% ✓
  2. Erro em regime estacionário nulo para entrada degrau: ess ≈ 0 ✓

O sistema possui realimentação negativa com integrador no controlador,
garantindo erro nulo em regime estacionário para entradas do tipo degrau.
=============================================================================
\end{lstlisting}
\subsection*{Gráficos}
\begin{figure}[H]
\centering
\includegraphics[width=0.95\linewidth]{figs/4\_5\_fig01.png}
\caption{PC4.5}
\end{figure}
\begin{figure}[H]
\centering
\includegraphics[width=0.95\linewidth]{figs/4\_5\_fig02.png}
\caption{PC4.5}
\end{figure}

\clearpage
\section*{PC4.8}
\noindent\textit{Fonte:} 4.8.ipynb
\vspace{0.4cm}
\subsection*{Resultados (saída do código)}
\begin{lstlisting}
=============================================================================

=============================================================================
PROBLEMA PC4.8 - PARTE (c): COMPARAÇÃO ENTRE CONTROLADORES
=============================================================================

TABELA COMPARATIVA:
-----------------------------------------------------------------------------
                      Aspecto Parte (a) - Proporcional     Parte (b) - PI
          Tipo de Controlador               G_c(s) = K G_c(s) = K₀ + K₁/s
                 Complexidade                  Simples      Mais complexo
         Número de Parâmetros                  1 (K=2)    2 (K₀=2, K₁=20)
              Ganho DC (y_ss)                 0.666666           0.999955
    Erro em Regime Est. (ess)                 0.333334         0.00004540
Erro em % do sinal de entrada                 33.3334%          0.004540%
   Tempo de resposta (aprox.)          Rápido (~0.03s)    Rápido (~0.02s)

ANÁLISE E DISCUSSÃO:
-----------------------------------------------------------------------------

1. DESEMPENHO DO ERRO EM REGIME ESTACIONÁRIO:
   • Controlador Proporcional (a):
     - Erro em regime estacionário: 0.333334 (33.3334%)
     - O sistema NÃO rastreia perfeitamente a entrada degrau
     - Há um erro residual permanente

   • Controlador PI (b):
     - Erro em regime estacionário: 0.00004540 (praticamente zero)
     - O sistema RASTREIA perfeitamente a entrada degrau
     - O integrador elimina qualquer erro estacionário

2. ANÁLISE TEÓRICA:
   • Controlador Proporcional:
     - Ganho em malha aberta finito: K*Gp(0) = 2*1 = 2
     - Ganho em malha fechada: T(0) = 2/(1+2) ≈ 0.667
     - Erro estacionário: 1 - 0.667 = 0.333 (por análise teórica)

   • Controlador PI:
     - Ganho em malha aberta infinito (devido ao integrador)
     - Ganho em malha fechada: T(0) = 1
     - Erro estacionário: 0 (eliminado pelo integrador)

3. TRADE-OFF COMPLEXIDADE vs DESEMPENHO:
   • Controlador Proporcional:
     ✓ Simplicidade: apenas 1 parâmetro (K)
     ✓ Implementação trivial
     ✗ Erro em regime estacionário não nulo
     ✗ Não atende ao requisito de rastreamento perfeito

   • Controlador PI:
     ✓ Rastreamento perfeito (ess = 0)
     ✓ Atende ao requisito do problema
     ✗ Complexidade aumentada: 2 parâmetros (K₀, K₁)
     ✗ Implementação requer integrador

4. CONCLUSÃO:
   O CONTROLADOR PI é a solução de COMPLEXIDADE MÍNIMA que permite
   rastreamento de referência com erro em regime estacionário NULO.

   • Aumento de complexidade: 1 parâmetro → 2 parâmetros
   • Benefício: Erro de 0.3333 → 0.00004540
   • Trade-off: A adição de um integrador (K₁/s) é essencial e mínima
     para eliminar o erro estacionário.

=============================================================================
\end{lstlisting}
\subsection*{Gráficos}
\begin{figure}[H]
\centering
\includegraphics[width=0.95\linewidth]{figs/4\_8\_fig01.png}
\caption{PC4.8}
\end{figure}
\begin{figure}[H]
\centering
\includegraphics[width=0.95\linewidth]{figs/4\_8\_fig02.png}
\caption{PC4.8}
\end{figure}
\begin{figure}[H]
\centering
\includegraphics[width=0.95\linewidth]{figs/4\_8\_fig03.png}
\caption{PC4.8}
\end{figure}

\clearpage
\section*{PC4.11}
\noindent\textit{Fonte:} 4.11.ipynb
\vspace{0.4cm}
\subsection*{Resultados (saída do código)}
\begin{lstlisting}
=====================================================================
PROBLEMA PC4.11 - ANÁLISE DO SISTEMA DE CONTROLE
=====================================================================

(a) FUNÇÃO DE TRANSFERÊNCIA EM MALHA FECHADA

Para K = 10:
T(s) = Y(s)/R(s) =
<TransferFunction>: sys[4]
Inputs (1): ['u[0]']
Outputs (1): ['y[0]']

          200 s + 200
  ----------------------------
  s^3 + 5.5 s^2 + 10.5 s + 206

Numerador: [200. 200.]
Denominador: [array([  1. ,   5.5,  10.5, 206. ])]

Para K = 12:
T(s) = Y(s)/R(s) =
<TransferFunction>: sys[9]
Inputs (1): ['u[0]']
Outputs (1): ['y[0]']

          240 s + 240
  ----------------------------
  s^3 + 5.5 s^2 + 10.5 s + 246

Numerador: [240. 240.]
Denominador: [array([  1. ,   5.5,  10.5, 246. ])]

Para K = 15:
T(s) = Y(s)/R(s) =
<TransferFunction>: sys[14]
Inputs (1): ['u[0]']
Outputs (1): ['y[0]']

          300 s + 300
  ----------------------------
  s^3 + 5.5 s^2 + 10.5 s + 306

Numerador: [300. 300.]
Denominador: [array([  1. ,   5.5,  10.5, 306. ])]

======================================================================
(b) RESPOSTA AO DEGRAU UNITÁRIO E ANÁLISE
======================================================================

─────────────────────────────────────────────────────────────────
K = 10
─────────────────────────────────────────────────────────────────
  Valor em regime estacionário (Y_ss):     -387.123196
  Erro em regime estacionário (e_ss):      388.123196
  Percentual de ultrapassagem (overshoot): 0.00%
  Tempo de acomodação (2%):                5.0000 s

─────────────────────────────────────────────────────────────────
K = 12
─────────────────────────────────────────────────────────────────
  Valor em regime estacionário (Y_ss):     2157.875664
  Erro em regime estacionário (e_ss):      -2156.875664
  Percentual de ultrapassagem (overshoot): 0.00%
  Tempo de acomodação (2%):                4.9900 s
  Tempo até 90% do valor final:            4.9800 s

─────────────────────────────────────────────────────────────────
K = 15
─────────────────────────────────────────────────────────────────
  Valor em regime estacionário (Y_ss):     -789.734395
  Erro em regime estacionário (e_ss):      790.734395
  Percentual de ultrapassagem (overshoot): 0.00%
  Tempo de acomodação (2%):                5.0000 s

======================================================================
Gráfico salvo como 'step_response_analysis.png'
======================================================================
\end{lstlisting}
\subsection*{Gráficos}
\begin{figure}[H]
\centering
\includegraphics[width=0.95\linewidth]{figs/4\_11\_fig01.png}
\caption{PC4.11}
\end{figure}

\clearpage
\section*{PC5.2}
\noindent\textit{Fonte:} 5.2.ipynb
\vspace{0.4cm}
\subsection*{Resultados (saída do código)}
\begin{lstlisting}
===========================================================================
PROBLEMA PC5.2 - RESPOSTA A RAMPA UNITÁRIA COM REALIMENTAÇÃO
===========================================================================

(a) FUNÇÃO DE TRANSFERÊNCIA EM MALHA ABERTA:
---------------------------------------------------------------------------
L(s) = G_c(s)G(s) =
<TransferFunction>: sys[0]
Inputs (1): ['u[0]']
Outputs (1): ['y[0]']

     s + 10
  ------------
  s^3 + 15 s^2

(b) FUNÇÃO DE TRANSFERÊNCIA EM MALHA FECHADA:
---------------------------------------------------------------------------
T(s) = L(s) / (1 + L(s)) = Y(s) / R(s) =
<TransferFunction>: sys[2]
Inputs (1): ['u[0]']
Outputs (1): ['y[0]']

         s + 10
  ---------------------
  s^3 + 15 s^2 + s + 10

Numerador: [ 1 10]
Denominador: [array([ 1, 15,  1, 10])]

(c) ANÁLISE DA RESPOSTA À RAMPA UNITÁRIA:
---------------------------------------------------------------------------

Intervalo de tempo: 0 ≤ t ≤ 50 s

No tempo final (t = 50.0 s):
  Entrada r(t) = 50.0000
  Saída y(t) = 50.0074
  Erro e(t) = r(t) - y(t) = -0.007395

(d) ERRO DE REGIME ESTACIONÁRIO TEÓRICO:
---------------------------------------------------------------------------

Constante de erro de velocidade (Velocity Error Constant):
K_v = lim(s→0) s*L(s) = inf

Erro de regime estacionário para rampa unitária:
e_ss = 1/K_v = 0.000000

Erro simulado ao final de 50 segundos: -0.007395
Diferença: 0.007395

===========================================================================
Gráficos salvos como 'rampa_response.png'
===========================================================================
\end{lstlisting}
\subsection*{Gráficos}
\begin{figure}[H]
\centering
\includegraphics[width=0.95\linewidth]{figs/5\_2\_fig01.png}
\caption{PC5.2}
\end{figure}

\clearpage
\section*{PC5.4}
\noindent\textit{Fonte:} 5.4.ipynb
\vspace{0.4cm}
\subsection*{Resultados (saída do código)}
\begin{lstlisting}
===========================================================================
PROBLEMA PC5.4 - ANÁLISE DO MÁXIMO PERCENTUAL DE ULTRAPASSAGEM
===========================================================================

(a) ANÁLISE ANALÍTICA:
---------------------------------------------------------------------------

Sistema em Malha Aberta:
  G_c(s) = 21/s (Controlador)
  G_p(s) = 1/(s+2) (Processo)
  L(s) = G_c(s)*G_p(s) =
<TransferFunction>: sys[0]
Inputs (1): ['u[0]']
Outputs (1): ['y[0]']

     21
  ---------
  s^2 + 2 s

Sistema em Malha Fechada:
  T(s) = L(s)/(1 + L(s)) = Y(s)/R(s) =
<TransferFunction>: sys[2]
Inputs (1): ['u[0]']
Outputs (1): ['y[0]']

        21
  --------------
  s^2 + 2 s + 21

  Numerador: [21]
  Denominador: [array([ 1,  2, 21])]

─────────────────────────────────────────────────────────────────────────────
PARÂMETROS DO SISTEMA DE SEGUNDA ORDEM:
─────────────────────────────────────────────────────────────────────────────
  Frequência Natural: ωn = sqrt(21) = 4.582576 rad/s
  Amortecimento Relativo: ζ = 1/sqrt(21) = 0.218218

  sqrt(1 - ζ²) = sqrt(1 - 0.047619) = 0.975900
  e^(-ζπ/sqrt(1-ζ²)) = e^(-0.218218*π/0.975900)
                   = e^(-0.702481)
                   = 0.495355

  MÁXIMA ULTRAPASSAGEM PERCENTUAL (M.U.P.):
  M.U.P. = 0.495355 * 100% = 49.54%

─────────────────────────────────────────────────────────────────────────────
OUTROS PARÂMETROS DO SISTEMA:
─────────────────────────────────────────────────────────────────────────────
  Frequência natural: 0.7293 Hz
  Frequência de oscilação: 4.4721 rad/s
  Tempo de pico (máxima ultrapassagem): t_p = 0.702481 s
  Tempo de acomodação (2%): t_s = 4.000000 s

=============================================================================
(b) RESPOSTA AO DEGRAU E ESTIMATIVA GRÁFICA
=============================================================================

A partir da simulação:
  Valor de pico: y_max = 1.495338
  Tempo de pico: t_p = 0.700701 s
  Ultrapassagem percentual: M.U.P. = 49.53%
  Tempo de acomodação (2%): t_s ≈ 4.319319 s

=============================================================================
COMPARAÇÃO ENTRE ANÁLISE TEÓRICA E SIMULAÇÃO:
=============================================================================

  Ultrapassagem percentual teórica (a):  49.54%
  Ultrapassagem percentual simulada (b): 49.53%
  Erro relativo:                         0.0017%

  Tempo de pico teórico:  0.702481 s
  Tempo de pico simulado: 0.700701 s

  Tempo de acomodação teórico (2%):  4.000000 s
  Tempo de acomodação simulado (2%): 4.319319 s

✓ Gráfico salvo como 'step_response_pc54.png'
=============================================================================
\end{lstlisting}
\subsection*{Gráficos}
\begin{figure}[H]
\centering
\includegraphics[width=0.95\linewidth]{figs/5\_4\_fig01.png}
\caption{PC5.4}
\end{figure}

\clearpage
\section*{PC5.8}
\noindent\textit{Fonte:} 5.8.ipynb
\vspace{0.4cm}
\subsection*{Resultados (saída do código)}
\begin{lstlisting}
=============================================================================
PROBLEMA PC5.8 - ANÁLISE DO PILOTO AUTOMÁTICO DE MÍSSIL
=============================================================================

(1) COMPONENTES DO SISTEMA:
-----------------------------------------------------------------------------

Controlador: G_c(s) = 0.1 + 5/s =
<TransferFunction>: sys[16]
Inputs (1): ['u[0]']
Outputs (1): ['y[0]']

  0.1 s + 5
  ---------
      s

Dinâmica do Míssil: G_p(s) = 100(s+1)/(s² + 2s + 100) =
<TransferFunction>: sys[17]
Inputs (1): ['u[0]']
Outputs (1): ['y[0]']

    100 s + 100
  ---------------
  s^2 + 2 s + 100

Função de Transferência em Malha Aberta:
L(s) = G_c(s) * G_p(s) =
<TransferFunction>: sys[18]
Inputs (1): ['u[0]']
Outputs (1): ['y[0]']

  10 s^2 + 510 s + 500
  --------------------
  s^3 + 2 s^2 + 100 s

Numerador de L(s): [ 10. 510. 500.]
Denominador de L(s): [array([  1,   2, 100,   0])]

=============================================================================
FUNÇÃO DE TRANSFERÊNCIA EM MALHA FECHADA:
=============================================================================
T(s) = L(s)/(1 + L(s)) = Theta_hat(s)/Theta_dot_d(s) =
<TransferFunction>: sys[20]
Inputs (1): ['u[0]']
Outputs (1): ['y[0]']

     10 s^2 + 510 s + 500
  --------------------------
  s^3 + 12 s^2 + 610 s + 500

Numerador de T(s): [ 10. 510. 500.]
Denominador de T(s): [  1.  12. 610. 500.]

=============================================================================
ANÁLISE DO SISTEMA:
=============================================================================

Pólos do sistema em malha fechada:
  Pólo 1: -5.583822 ± j23.864733
  Pólo 2: -5.583822 ± j23.864733
  Pólo 3: -0.832356

Pólos dominantes (conjugado complexo):
  s = --5.583822 ± j23.864733

Parâmetros equivalentes de segunda ordem (pólos dominantes):
  ωn (frequência natural estimada) = 24.509275 rad/s
  ζ (amortecimento relativo estimado) = 0.227825

=============================================================================
PREDIÇÃO USANDO FÓRMULAS DE SEGUNDA ORDEM:
=============================================================================

Parâmetros teóricos de segunda ordem:
  Máximo percentual de ultrapassagem:
    M_p = e^(-ζπ/sqrt(1-ζ²)) = e^(-0.227825*π/0.973702)
        = e^(-0.735064)
        = 0.479475
    M_p = 47.95%

  Tempo até o pico:
    T_p = π/(ωn*sqrt(1-ζ²)) = π/(24.509275*0.973702)
        = 0.131642 s

  Tempo de acomodação (2%):
    T_s = 4/(ζ*ωn) = 4/(0.227825*24.509275)
        = 0.716355 s

=============================================================================
(b) RESPOSTA AO DEGRAU UNITÁRIO - ANÁLISE REAL:
=============================================================================

RESULTADOS DA SIMULAÇÃO:
  Valor final (regime estacionário): 0.986231
  Valor de pico: 1.297871
  Tempo de pico: 0.111056 s
  Percentual de ultrapassagem (real): M_p = 29.79%
  Tempo de acomodação (2%, real): T_s ≈ 1.931466 s

=============================================================================
COMPARAÇÃO - PREDIÇÃO (2ª ORDEM) vs RESPOSTA REAL:
=============================================================================

Parâmetro
...
(continuação omitida por brevidade)
\end{lstlisting}
\subsection*{Gráficos}
\begin{figure}[H]
\centering
\includegraphics[width=0.95\linewidth]{figs/5\_8\_fig01.png}
\caption{PC5.8}
\end{figure}

\clearpage
\section*{PC5.9}
\noindent\textit{Fonte:} 5.9.ipynb
\vspace{0.4cm}
\subsection*{Resultados (saída do código)}
\begin{lstlisting}
=============================================================================
PROBLEMA PC5.9 - ANÁLISE DE SISTEMA COM REALIMENTAÇÃO NÃO-UNITÁRIA
=============================================================================

(1) COMPONENTES DO SISTEMA:
-----------------------------------------------------------------------------

Caminho direto: G(s) = 10/(s+10) =
<TransferFunction>: sys[0]
Inputs (1): ['u[0]']
Outputs (1): ['y[0]']

    10
  ------
  s + 10

Caminho de realimentação: H(s) = 0.5/(10s+0.5) =
<TransferFunction>: sys[1]
Inputs (1): ['u[0]']
Outputs (1): ['y[0]']

     0.5
  ----------
  10 s + 0.5

H(s) simplificado: H(s) = 0.05/(s+0.05) =
<TransferFunction>: sys[2]
Inputs (1): ['u[0]']
Outputs (1): ['y[0]']

    0.05
  --------
  s + 0.05

=============================================================================
FUNÇÃO DE TRANSFERÊNCIA EM MALHA ABERTA:
=============================================================================
L(s) = G(s)*H(s) =
<TransferFunction>: sys[3]
Inputs (1): ['u[0]']
Outputs (1): ['y[0]']

           5
  --------------------
  10 s^2 + 100.5 s + 5

Numerador de L(s): [5.]
Denominador de L(s): [ 10.  100.5   5. ]

=============================================================================
FUNÇÃO DE TRANSFERÊNCIA EM MALHA FECHADA:
=============================================================================
T(s) = G(s)/(1 + G(s)*H(s)) = Y(s)/R(s) =
<TransferFunction>: sys[4]
Inputs (1): ['u[0]']
Outputs (1): ['y[0]']

        100 s + 5
  ---------------------
  10 s^2 + 100.5 s + 10

Numerador de T(s): [100.   5.]
Denominador de T(s): [ 10.  100.5  10. ]

=============================================================================
ANÁLISE DO SISTEMA:
=============================================================================

Pólos do sistema em malha fechada:
  Pólo 1: -9.949492
  Pólo 2: -0.100508

Estabilidade: ESTÁVEL

=============================================================================
RESPOSTA AO DEGRAU UNITÁRIO:
=============================================================================

RESULTADOS:
  Valor final (regime estacionário): 0.686753
  Valor de pico: 0.978591
  Tempo de pico: 0.535268 s
  Percentual de ultrapassagem (M.U.P.): 42.50%
  Tempo de acomodação (2%): 9.289645 s
  Tempo de acomodação (5%): 8.319160 s

=============================================================================
CARACTERÍSTICAS DO SISTEMA:
=============================================================================

Sistema é de 2ª ordem com dois pólos reais (superamortecido)

=============================================================================
RESUMO DOS RESULTADOS:
=============================================================================

✓ Máxima Ultrapassagem Percentual (M.U.P.): 42.50%
✓ Tempo de Acomodação (Critério 2%):    9.289645 s
✓ Tempo de Acomodação (Critério 5%):    8.319160 s
✓ Tempo de Pico:                         0.535268 s
✓ Sistema é ESTÁVEL

✓ Gráfico salvo como 'system_pc59_response.png'
=============================================================================
\end{lstlisting}
\subsection*{Gráficos}
\begin{figure}[H]
\centering
\includegraphics[width=0.95\linewidth]{figs/5\_9\_fig01.png}
\caption{PC5.9}
\end{figure}

\clearpage
\section*{PC5.10}
\noindent\textit{Fonte:} 5.10.ipynb
\vspace{0.4cm}
\subsection*{Resultados (saída do código)}
\begin{lstlisting}
=============================================================================
PROBLEMA PC5.10 - RESPOSTA A RAMPA E ERRO EM REGIME ESTACIONÁRIO
=============================================================================

(1) FUNÇÃO DE TRANSFERÊNCIA EM MALHA ABERTA:
-----------------------------------------------------------------------------

L(s) = G(s) = 10/(s(s+15)(s+5)) =
<TransferFunction>: sys[7]
Inputs (1): ['u[0]']
Outputs (1): ['y[0]']

          10
  -------------------
  s^3 + 20 s^2 + 75 s

=============================================================================
FUNÇÃO DE TRANSFERÊNCIA EM MALHA FECHADA:
=============================================================================
T(s) = L(s)/(1 + L(s)) = Y(s)/R(s) =
<TransferFunction>: sys[9]
Inputs (1): ['u[0]']
Outputs (1): ['y[0]']

             10
  ------------------------
  s^3 + 20 s^2 + 75 s + 10

Numerador de T(s): [10]
Denominador de T(s): [ 1 20 75 10]

=============================================================================
ANÁLISE DE ESTABILIDADE:
=============================================================================

Pólos do sistema em malha fechada:
  Pólo 1: -15.065940
  Pólo 2: -4.795654
  Pólo 3: -0.138406

Estabilidade: ESTÁVEL

=============================================================================
ERRO EM REGIME ESTACIONÁRIO:
=============================================================================

Constante de erro de velocidade:
  K_v = lim(s→0) s*L(s) = 10/(15*5) = 0.133333

Erro em regime estacionário para rampa unitária:
  e_ss = 1/K_v = 1/0.133333 = 7.500000

=============================================================================
SIMULAÇÃO DA RESPOSTA À RAMPA:
=============================================================================

RESULTADOS DA SIMULAÇÃO (t = 100.0 s):
  Entrada r(t) = 100.000000
  Saída y(t) = 92.500007
  Erro e(t) = 7.499993

COMPARAÇÃO:
  Erro teórico:   7.500000
  Erro simulado:  7.499993
  Diferença:      0.000007

Tempo para atingir regime estacionário: ~100.00 s

=============================================================================
TIPO DE SISTEMA:
=============================================================================

Tipo de Sistema: Tipo 1 (uma integração)
  - Erro zero para entrada em degrau: SIM
  - Erro finito para entrada em rampa: SIM (e_ss = 1/K_v)
  - Erro infinito para entrada em parábola: SIM

=============================================================================
RESUMO FINAL:
=============================================================================

✓ Entrada: r(t) = t (rampa unitária)
✓ Erro em regime estacionário teórico: e_ss = 7.500000
✓ Erro em regime estacionário simulado: e_ss ≈ 7.499993
✓ Sistema: ESTÁVEL (Tipo 1)

✓ Gráficos salvos como 'ramp_response_pc510.png'
=============================================================================
\end{lstlisting}
\subsection*{Gráficos}
\begin{figure}[H]
\centering
\includegraphics[width=0.95\linewidth]{figs/5\_10\_fig01.png}
\caption{PC5.10}
\end{figure}

\clearpage
\section*{PC5.11}
\noindent\textit{Fonte:} 5.11.ipynb
\vspace{0.4cm}
\subsection*{Resultados (saída do código)}
\begin{lstlisting}
============================================================
PARTE (a): FUNÇÃO DE TRANSFERÊNCIA EM MALHA FECHADA
============================================================

Controlador G_c(s) = 0.5 + 2/s:
<TransferFunction>: sys[51]
Inputs (1): ['u[0]']
Outputs (1): ['y[0]']

  0.5 s + 2
  ---------
      s

Processo G_p(s) = 1/(s(s+2)):
<TransferFunction>: sys[56]
Inputs (1): ['u[0]']
Outputs (1): ['y[0]']

      1
  ---------
  s^2 + 2 s

Função de transferência em malha aberta G_ol(s) = G_c(s)*G_p(s):
<TransferFunction>: sys[57]
Inputs (1): ['u[0]']
Outputs (1): ['y[0]']

   0.5 s + 2
  -----------
  s^3 + 2 s^2

Função de transferência em malha fechada T(s) = Y(s)/R(s):
<TransferFunction>: sys[59]
Inputs (1): ['u[0]']
Outputs (1): ['y[0]']

         0.5 s + 2
  -----------------------
  s^3 + 2 s^2 + 0.5 s + 2

Numerador: 0.5s + 2.0
Denominador: s³ + 2.0s² + 0.5s + 2.0

============================================================
PARTE (b): RESPOSTA A DIFERENTES ENTRADAS
============================================================

============================================================
ANÁLISE DE ESTABILIDADE
============================================================

Zeros do sistema:
  Zero 1: -4.000000

Polos do sistema em malha fechada:
  Polo 1: -2.188978
  Polo 2: 0.094489+0.951178j
  Polo 3: 0.094489-0.951178j

Sistema é estável? NÃO ✗
\end{lstlisting}
\subsection*{Gráficos}
\begin{figure}[H]
\centering
\includegraphics[width=0.95\linewidth]{figs/5\_11\_fig01.png}
\caption{PC5.11}
\end{figure}

\clearpage
\section*{PC5.12}
\noindent\textit{Fonte:} 5.12.ipynb
\vspace{0.4cm}
\subsection*{Resultados (saída do código)}
\begin{lstlisting}
======================================================================
PC5.12: ANÁLISE DE RESPOSTA AO DEGRAU E DESEMPENHO DO SISTEMA
======================================================================

======================================================================
PARTE (a): SISTEMA ORIGINAL
======================================================================

Função de transferência em malha fechada:
T(s) = 77(s + 2) / [(s + 7)(s² + 4s + 22)]

Sistema original:
<TransferFunction>: sys[15]
Inputs (1): ['u[0]']
Outputs (1): ['y[0]']

         77 s + 154
  -------------------------
  s^3 + 11 s^2 + 50 s + 154

Denominador expandido:
(s + 7)(s² + 4s + 22) = s³ + 11s² + 50s + 154

----------------------------------------------------------------------
CARACTERÍSTICAS DA RESPOSTA AO DEGRAU:
----------------------------------------------------------------------
Valor final em regime permanente: 1.039955
Máximo valor atingido: 1.806442
Máxima Ultrapassagem Percentual (MUP): 73.70%
Tempo de acomodação Ts (critério 2%): 1.8199 segundos

Polos do sistema original:
  Polo 1: -7.000000
  Polo 2: -2.000000+4.242641j
  Polo 3: -2.000000-4.242641j

======================================================================
PARTE (b): SISTEMA SIMPLIFICADO (Desprezando polo em s = -7)
======================================================================

Função de transferência simplificada:
T_b(s) = 77(s + 2) / (s² + 4s + 22)

Sistema simplificado:
<TransferFunction>: sys[40]
Inputs (1): ['u[0]']
Outputs (1): ['y[0]']

    77 s + 154
  --------------
  s^2 + 4 s + 22

----------------------------------------------------------------------
CARACTERÍSTICAS DA RESPOSTA AO DEGRAU:
----------------------------------------------------------------------
Valor final em regime permanente: 7.295222
Máximo valor atingido: 14.081377
Máxima Ultrapassagem Percentual (MUP): 93.02%
Tempo de acomodação Ts (critério 2%): 1.6418 segundos

Polos do sistema simplificado:
  Polo 1: -2.000000+4.242641j
  Polo 2: -2.000000-4.242641j

Análise de 2ª ordem:
Frequência natural ωn = 4.6904 rad/s
Coeficiente de amortecimento ζ = 0.4264
Tempo de acomodação teórico (2%): 2.0000 segundos
Máxima Ultrapassagem Teórica: 22.74%

======================================================================
COMPARAÇÃO E CONCLUSÕES
======================================================================

Parâmetro                                     | Sistema Original | Sistema Simplificado
-----------------------------------------------------------------------------
Tempo de acomodação Ts (s)                    | 1.8199          | 1.6418         
Máxima Ultrapassagem (MUP)                    | 73.70          % | 93.02          %
Valor final em regime permanente              | 1.039955        | 7.295222       

----------------------------------------------------------------------
DIFERENÇAS RELATIVAS:
----------------------------------------------------------------------
Diferença em Ts: 0.1781 s (9.79%)
Diferença em MUP: 19.32% (26.21%)

----------------------------------------------------------------------
CONCLUSÕES SOBRE O DESPREZAMENTO DO POLO EM s = -7:
----------------------------------------------------------------------

1. O polo em s = -7 está mais à esquerda no plano complexo que os p
...
(continuação omitida por brevidade)
\end{lstlisting}

\clearpage
\section*{PC6.2}
\noindent\textit{Fonte:} 6.2.ipynb
\vspace{0.4cm}
\subsection*{Resultados (saída do código)}
\begin{lstlisting}
=============================================================================
PC6.2: ANÁLISE DE ESTABILIDADE - RAÍZES DO POLINÔMIO CARACTERÍSTICO
=============================================================================

Função da Planta:
G(s) = (s² - 2s + 4)/(s² + 4s + 2)
<TransferFunction>: sys[17]
Inputs (1): ['u[0]']
Outputs (1): ['y[0]']

  s^2 - 2 s + 4
  -------------
  s^2 + 4 s + 2

=============================================================================
ANÁLISE DAS RAÍZES DO POLINÔMIO CARACTERÍSTICO
=============================================================================

-----------------------------------------------------------------------------
CASO: K = 1
-----------------------------------------------------------------------------

Controlador: Gc(s) = 1

Função de transferência em malha fechada:
<TransferFunction>: sys[21]
Inputs (1): ['u[0]']
Outputs (1): ['y[0]']

   s^2 - 2 s + 4
  ---------------
  2 s^2 + 2 s + 6

Polinômio característico: (2)s² + (2)s + (6) = 0
Simplificado: 2s² + 2s + 6 = 0

Raízes (Polos em malha fechada):
  Polo 1: -0.500000+1.658312j
  Polo 2: -0.500000-1.658312j

Análise de Estabilidade:
  Número de polos com parte real negativa: 2 de 2
  Sistema ESTÁVEL? SIM ✓

-----------------------------------------------------------------------------
CASO: K = 2
-----------------------------------------------------------------------------

Controlador: Gc(s) = 2

Função de transferência em malha fechada:
<TransferFunction>: sys[25]
Inputs (1): ['u[0]']
Outputs (1): ['y[0]']

  2 s^2 - 4 s + 8
  ---------------
    3 s^2 + 10

Polinômio característico: (3)s² + (0)s + (10) = 0
Simplificado: 3s² + 0s + 10 = 0

Raízes (Polos em malha fechada):
  Polo 1: -0.000000+1.825742j
  Polo 2: 0.000000-1.825742j

Análise de Estabilidade:
  Número de polos com parte real negativa: 0 de 2
  Sistema ESTÁVEL? NÃO ✗

-----------------------------------------------------------------------------
CASO: K = 5
-----------------------------------------------------------------------------

Controlador: Gc(s) = 5

Função de transferência em malha fechada:
<TransferFunction>: sys[29]
Inputs (1): ['u[0]']
Outputs (1): ['y[0]']

  5 s^2 - 10 s + 20
  -----------------
  6 s^2 - 6 s + 22

Polinômio característico: (6)s² + (-6)s + (22) = 0
Simplificado: 6s² + -6s + 22 = 0

Raízes (Polos em malha fechada):
  Polo 1: 0.500000+1.848423j
  Polo 2: 0.500000-1.848423j

Análise de Estabilidade:
  Número de polos com parte real negativa: 0 de 2
  Sistema ESTÁVEL? NÃO ✗

=============================================================================
RESUMO COMPARATIVO
=============================================================================

K          | Polinômio       | Raízes                         | Estável?    
-----------------------------------------------------------------------------
1          | 2s²+2s+6        | -0.500+1.658j, -0.500-1.658j   | ✓ SIM       
2          | 3s²+0s+10       | -0.000+1.826j, 0.000-1.826j    | ✗ NÃO       
5          | 6s²+-6s+22      | 0.500+1.848j, 0.500-1.848j     | ✗ NÃO       

=============================================================================
CONCLUSÕES SOBRE ESTABILIDADE
==================================================================
...
(continuação omitida por brevidade)
\end{lstlisting}
\subsection*{Gráficos}
\begin{figure}[H]
\centering
\includegraphics[width=0.95\linewidth]{figs/6\_2\_fig01.png}
\caption{PC6.2}
\end{figure}

\clearpage
\section*{PC6.4}
\noindent\textit{Fonte:} 6.4.ipynb
\vspace{0.4cm}
\subsection*{Resultados (saída do código)}
\begin{lstlisting}
Linha s^3 é toda zero! Usando polinômio auxiliar.
Linha s^1 é toda zero! Usando polinômio auxiliar.

Tabela de Routh-Hurwitz:
[[1. 2. 1.]
 [2. 4. 2.]
 [8. 8. 0.]
 [2. 2. 0.]
 [4. 0. 0.]
 [2. 0. 0.]]

Mudanças de sinal na 1ª coluna: 0
Polos do sistema:
[-2.00000000e+00+0.j         -3.97581489e-09+1.00000002j
 -3.97581489e-09-1.00000002j  3.97581464e-09+0.99999998j
  3.97581464e-09-0.99999998j]
Linha s^3 é toda zero! Usando polinômio auxiliar.
Linha s^1 é toda zero! Usando polinômio auxiliar.
\end{lstlisting}
\subsection*{Observação}
Como há uma \textbf{linha de zeros} na tabela de Routh e aparecem polos \textbf{no eixo imaginário} (numericamente, com parte real muito próxima de zero), o sistema não é assintoticamente estável: ele está em \textbf{estabilidade marginal} (oscilação sustentada).
\subsection*{Gráficos}
\begin{figure}[H]
\centering
\includegraphics[width=0.95\linewidth]{figs/6\_4\_fig01.png}
\caption{PC6.4}
\end{figure}

\clearpage
\section*{PC6.6}
\noindent\textit{Fonte:} 6.6.ipynb
\vspace{0.4cm}
\subsection*{Resultados (saída do código)}
\begin{lstlisting}
--- Análise para o limite de estabilidade (K = 4) ---
A equação característica de malha fechada é:
[1 5 1 5]

As raízes (polos) para K = 4 são:
Raiz 1: (-5+0j)
Raiz 2: 1j
Raiz 3: -1j
\end{lstlisting}
\subsection*{Gráficos}
\begin{figure}[H]
\centering
\includegraphics[width=0.95\linewidth]{figs/6\_6\_fig01.png}
\caption{PC6.6}
\end{figure}

\clearpage
\section*{PC6.8}
\noindent\textit{Fonte:} 6.8.ipynb
\vspace{0.4cm}
\subsection*{Resultados (saída do código)}
\begin{lstlisting}
ITEM (a) - Critério de Routh-Hurwitz
------------------------------------
Polinômio característico:
s^3 + 10 s^2 + 10 s + 5 K1 = 0

Tabela de Routh:
s^3 | 1        10
s^2 | 10       5K1
s^1 | (100-5K1)/10   0
s^0 | 5K1

Para estabilidade, todos elementos da 1ª coluna > 0:
1 > 0
10 > 0
(100 - 5K1)/10 > 0
5K1 > 0

Logo:
K1 > 0
K1 < 20

Faixa estável:
0 < K1 < 20

ITEM (b) - Comentários
----------------------
• Para 0 < K1 < 20: todos polos no semiplano esquerdo => sistema estável.
• Em K1 = 20: polos cruzam eixo imaginário => estabilidade marginal.
• Para K1 > 20: aparecem polos no semiplano direito => sistema instável.
• Ao aumentar K1, o sistema tende a ficar mais rápido inicialmente,
  porém perde estabilidade após K1 = 20.
\end{lstlisting}
\subsection*{Gráficos}
\begin{figure}[H]
\centering
\includegraphics[width=0.95\linewidth]{figs/6\_8\_fig01.png}
\caption{PC6.8}
\end{figure}

\end{document}
